% ========================================================================
% فصل ۱۳: مسائل باز و کارهای آینده
% منابع: farsimadanReviewSecurityChallenges2025، rahmawatiagustinaSystematicLiteratureReview2025،
% ghovanlooyghajarSBTMSScalableBlockchain2021، hozouriOverviewVANETVehicular2023،
% VanetSecurityPrivacy، al-aniPrivacySafetyImprovement2023، SecuringVehicularAd،
% mariaBAIVEfficientBlockchainBased2022، mazzocca2025survey
% ========================================================================
\section{مسائل باز و کارهای آینده}

بر اساس بررسی‌های انجام‌شده، مهم‌ترین مسائل باز در امنیت شبکه اقتضایی خودرویی و جهت‌گیری‌های پژوهشی برای رفع آن‌ها در حوزه‌های زیر دسته‌بندی می‌شوند \cite{farsimadanReviewSecurityChallenges2025, rahmawatiagustinaSystematicLiteratureReview2025}. این فصل مسائلی را گزارش می‌کند که منابع بررسی‌شده آن‌ها را به‌صورت صریح یا ضمنی مطرح کرده‌اند؛ در جایی که پوشش منابع ناکافی بوده، آن نکته در متن شفاف اعلام شده است.

\subsection{معاوضه حریم خصوصی، امنیت و کارایی}

چالش مرکزی در طراحی سامانه‌های امنیتی شبکه اقتضایی خودرویی، برقراری تعادل میان اهداف متعارض است. مرور ره‌آواتی آگوستینا و همکاران هنوز «طراحی طرحی که سه مسئلهٔ کلیدی حریم خصوصی، ایمنی و کارایی را به‌طور مؤثر برآورده کند» را چالشی علمی می‌نامد و نشان می‌دهد افزایش گمنامی معمولاً به قیمت کاهش اشتراک‌گذاری وضعیت در دوره‌های سکوت تمام می‌شود \cite{al-aniPrivacySafetyImprovement2023}. از سوی دیگر، الزام دسترس‌پذیری استفاده از تکنیک‌های رمزنگاری سبک و کم‌هزینه را ضروری می‌کند و سربار محاسباتی مکانیزم‌های امنیتی سنگین، اجرای آن‌ها بر خودروها را با مشکل مواجه می‌کند \cite{farsimadanReviewSecurityChallenges2025}. تعادل میان کاهش مؤثر تهدیدات و حداقل تأثیر بر عملکرد سامانه و راحتی کاربر، از دغدغه‌های اصلی معماری‌های امنیتی است \cite{SecuringVehicularAd}.

\subsection{استراتژی تغییر شبه‌نام و چرخه حیات کامل آن}

استراتژی‌های تغییر شبه‌نام در منابع بررسی‌شده با محدودیت‌هایی همراه‌اند. راهبردهایی که با بازهٔ زمانی ثابت کار می‌کنند و چگالی ترافیک را نادیده می‌گیرند، به حملهٔ پیوند منجر می‌شوند و چون ورودی‌های تولید شبه‌نام ثابت‌اند، پیوند دادن شبه‌نام‌ها — به‌ویژه اگر بازهٔ تغییر شبه‌نام قابل پیش‌بینی باشد — آسان می‌گردد \cite{rahmawatiagustinaSystematicLiteratureReview2025}. مرور مذکور افزودن ورودی‌های پویا، مانند شرایط جاده، مسیر خودروها و موقعیت‌های جغرافیایی را پیشنهاد می‌کند و در همان حال هشدار می‌دهد که باید میان غیرقابل‌پیش‌بینی بودن و ردیابی‌پذیری تعادل برقرار شود \cite{rahmawatiagustinaSystematicLiteratureReview2025}. از این رو، پیاده‌سازی کامل چرخهٔ حیات شبه‌نام — صدور، استفاده، تغییر، تفکیک و ابطال — بدون افت عملکرد سامانه، هدف کلیدی پژوهش‌های آینده است \cite{rahmawatiagustinaSystematicLiteratureReview2025}. این مسئله شامل یکپارچه‌سازی استراتژی‌های تغییر شبه‌نام با سایر مکانیزم‌های امنیتی نیز هست؛ نبود مرحلهٔ لغو در چرخهٔ حیات، خطر حملهٔ سیبل و دسترسی غیرمجاز را بالا می‌برد و نبود مرحلهٔ تفکیک، مسئلهٔ مرجع حل هویت را ایجاد می‌کند \cite{rahmawatiagustinaSystematicLiteratureReview2025}.

\subsection{اتکا به واحدهای کنارجاده‌ای و نقطه شکست منفرد}

وابستگی به واحدهای کنارجاده‌ای \LTRfootnote{\lr{Road-Side Unit (RSU)}} و مراجع متمرکز به‌عنوان گره‌های مورد اعتماد، نقطهٔ شکست منفرد \LTRfootnote{\lr{Single Point of Failure (SPoF)}} ایجاد می‌کند. در سامانه‌های مدیریت اعتماد متمرکز، گلوگاه و نقطهٔ شکست منفرد وجود دارد و اگر سامانهٔ اصلی از کار بیفتد، کل سامانه خراب می‌شود \cite{ghovanlooyghajarSBTMSScalableBlockchain2021}. در مدل استاندارد، عدم در دسترس بودن مرجع قابل اعتماد بر همهٔ موجودیت‌ها اثر می‌گذارد و چون این مرجع هویت همهٔ نهادها را می‌داند، هدف اصلی نقض حریم خصوصی می‌شود و تسخیر آن می‌تواند نقض گستردهٔ حریم خصوصی به همراه داشته باشد \cite{rahmawatiagustinaSystematicLiteratureReview2025}. از منظر شبکه، کانال‌های پهنای باند و واحدهای کنارجاده‌ای منابعی مشترک هستند و استقرار شبکه اقتضایی خودرویی در پی این اشتراک‌گذاری با چالش‌هایی روبه‌رو می‌شود \cite{hozouriOverviewVANETVehicular2023}. راه‌حل بررسی‌شده، تفکیک مرجع به مرجع ردیابی — که تولید شبه‌هویت و ردیابی خودروهای مخرب را بر عهده دارد — و مرکز تولید کلید — که تولید کلیدهای رمزنگاری و پارامترهای امنیتی را انجام می‌دهد — است تا هیچ نهاد واحدی بتواند هم هویت کاربران را فاش کند و هم کلیدهای رمزنگاری را در اختیار داشته باشد \cite{rahmawatiagustinaSystematicLiteratureReview2025}.

\subsection{مقیاس‌پذیری، محدودیت منابع و معماری‌های ترکیبی}

مقیاس‌پذیری پروتکل‌های امنیتی در ترافیک شهری متراکم به چالش کشیده می‌شود. توان محاسباتی خودرو نسبت به رایانه، تبلت و گوشی محدودتر است و برخی سازوکارهای امنیتی به‌دلیل سربار بالا قابل اجرا نیستند \cite{farsimadanReviewSecurityChallenges2025}. در حوزهٔ دفتر کل توزیع‌شده، بیشتر فناوری‌های زنجیره بلوکی به‌دلیل فرآیند استخراج، کاستی‌هایی مانند سرپایی پایین تراکنش و مقیاس‌پذیری ضعیف دارند، در حالی که در شبکه اقتضایی خودرویی تعداد و سرعت تولید تراکنش‌ها بسیار زیاد است \cite{ghovanlooyghajarSBTMSScalableBlockchain2021}. الگوریتم شاردینگ با تقسیم وظایف استخراج به کمیته‌ها، مقیاس‌پذیری را به‌صورت تقریباً خطی به توان محاسباتی واحدهای کنارجاده‌ای ارتقا می‌دهد \cite{ghovanlooyghajarSBTMSScalableBlockchain2021}. در این راستا، گسترش معماری‌های ترکیبی ابر، مه و زنجیره بلوک برای تأمین هم‌زمان مقیاس‌پذیری و امنیت، مسیری پژوهشی است؛ مرور ره‌آواتی آگوستینا و همکاران مدل ترکیبی «ابر + مه + بلاکچین» را با چالش‌های پیچیدگی، هزینه و تعامل‌پذیری گزارش می‌کنند \cite{rahmawatiagustinaSystematicLiteratureReview2025}.

\subsection{واقع‌گرایی ارزیابی‌ها و مجموعه‌داده‌های باز}

بسیاری از ارزیابی‌های امنیتی در محیط‌های شبیه‌سازی‌شده با فرضیات ساده‌شده انجام می‌شوند. مرور فارسیمدان و همکاران گزارش می‌کنند که اکثر روش‌ها با شبیه‌سازی ارزیابی می‌شوند، شبیه‌سازها ضیف‌اند و داده‌ها به‌صورت عمومی در دسترس قرار نمی‌گیرند \cite{farsimadanReviewSecurityChallenges2025}. حتی در طرح‌هایی که مجموعه‌داده استاندارد به کار برده‌اند — مانند مجموعه‌دادهٔ \lr{V2V} اتحادیهٔ اروپایی در سامانهٔ \lr{DIVA} — ارزیابی محدود به سناریوهای شبیه‌سازی‌شده و پیام‌های اعلان محیطی بوده و دقت‌ها تک‌سناریویی‌اند \cite{feraudoDIVADIDbasedReputation2024}. در برخی طرح‌ها نیز بلاکچین در ارزیابی عملی پیاده‌سازی اصلاً به کار گرفته نشده و تنها عملیات رمزنگاری زمان‌سنجی شده است \cite{mariaBAIVEfficientBlockchainBased2022}. بنابراین ایجاد و انتشار مجموعه‌داده‌های استاندارد و واقعی و معیارهای ارزیابی یکپارچه برای اعتبارسنجی روش‌های امنیتی، از نیازهای پژوهشی است \cite{farsimadanReviewSecurityChallenges2025}.

\subsection{مقاومت مدل‌های یادگیری ماشین و یادگیری فدرال}

مدل‌های یادگیری ماشین به کار رفته در تشخیص ناهنجاری و تشخیص نفوذ، در برابر حملات تخاصمی و مسموم‌سازی داده \LTRfootnote{\lr{Data Poisoning}} آسیب‌پذیر هستند. مرور فارسیمدان و همکاران گزارش می‌کنند که یادگیری فدرال به ساختار کارگزار–کارخواه وابسته است و راهبردهای تخاصمی و مسموم‌سازی داده بر مدل‌های محلی اثر می‌گذارند؛ جهت آیندهٔ این حوزه، طراحی مدل‌های مقاوم و حفظ‌کنندهٔ حریم خصوصی اعلام شده است \cite{farsimadanReviewSecurityChallenges2025}. در همین حال، شبکه‌های مولد تخاخمی \LTRfootnote{\lr{Generative Adversarial Network (GAN)}} و یادگیری تقویتی نیز در تشخیص رفتارهای غیرعادی و ارزیابی اعتمادپذیری خودروهای خودران و شناسایی نفوذگران با یادگیری \lr{Q} کاربرد داشته‌اند \cite{farsimadanReviewSecurityChallenges2025}. باید توجه داشت که جزئیات بیشتر دربارهٔ مقاوم‌سازی یادگیری فدرال در منابع بررسی‌شده پوشش داده نشده و این شکاف در خود منبع هم اعلام شده است \cite{farsimadanReviewSecurityChallenges2025}.

\subsection{آسیب‌پذیری‌های بی‌سیم و امنیت لایهٔ مسیریابی}

استرقاق سمع \LTRfootnote{\lr{Eavesdropping}}، حملهٔ فردی در میان \LTRfootnote{\lr{Man-in-the-Middle (MiM)}} و حملهٔ دوقلوی شیطانی \LTRfootnote{\lr{Evil Twin}} از جمله تهدیدات ذاتی کانال بی‌سیم هستند و در شبکه اقتضایی خودرویی که ارتباطات با برد کوتاه و در محیط‌های متحرک انجام می‌شود، مقابله با این حملات دشوارتر است \cite{farsimadanReviewSecurityChallenges2025}. در لایهٔ مسیریابی، مینتمور و سن چهار حملهٔ گودال سیاه، رهاکردن بسته، سیل و اطلاعات نادرست را بر روی پروتکل‌های مسیریابی \lr{AODV} و \lr{GPSR} شبیه‌سازی کرده‌اند \cite{VanetSecurityPrivacy} و مسیریابی داده یکی از سه مسئلهٔ بنیادین ارتباط میان‌خودرویی به شمار می‌رود \cite{hozouriOverviewVANETVehicular2023}. از این رو، توسعهٔ مکانیزم‌های مقاوم برای امنیت ارتباطات چندپرشی در لایه‌های مختلف سامانه‌های حمل‌ونقل هوشمند، از نیازهای پژوهشی است؛ این نیاز از تهدیدات گزارش‌شده در منابع استنباط می‌شود، هرچند پوشش مفصل سازوکارهای مقاوم چندپرشی در منابع بررسی‌شده محدود است \cite{farsimadanReviewSecurityChallenges2025, VanetSecurityPrivacy}.

\subsection{امضاهای دیجیتال سریع و سبک}

طرح‌های امضای دیجیتال موجود یا سرعت کافی برای کاربردهای بلادرنگ ایمنی را ندارند یا از نظر محاسباتی برای واحدهای سوارشونده با منابع محدود سنگین هستند. الزام دسترس‌پذیری به‌صراحت استفاده از تکنیک‌های رمزنگاری سبک و کم‌هزینه را ضروری می‌کند \cite{farsimadanReviewSecurityChallenges2025} و بهینه‌سازی تأیید امضاها برای کاهش سربار محاسباتی در بررسی پیام‌های ایمنی پیشنهاد شده است \cite{SecuringVehicularAd}. در طرح‌های مبتنی بر شبه‌نام نیز هنگام تغییر شبه‌نام یا آشنایی با همسایهٔ جدید، گواهی باید پیوست و تأیید شود و این امر سربار امنیتی را افزایش می‌دهد \cite{al-aniPrivacySafetyImprovement2023}. طرح‌های مبتنی بر رمزنگاری خم بیضوی، هزینهٔ احراز هویت یک کاربر و یک واحد کنارجاده‌ای را به حدود ۱۲.۵۴ میلی‌ثانیه — تقریباً ۸۰ کاربر در ثانیه — پایین آورده‌اند \cite{mariaBAIVEfficientBlockchainBased2022}. طراحی و استانداردسازی امضاهای سبک و سریع که امنیت کافی را نیز حفظ کنند، همچنان یک نیاز برآورده‌نشده است؛ باید توجه داشت که ادعاهای دربارهٔ امضای تجمیعی و امضای بدون گواهی در منابع بررسی‌شده پشتیبانی نشده و در فصل ۶ به همین شکاف اشاره شد \cite{farsimadanReviewSecurityChallenges2025, SecuringVehicularAd}.

\subsection{تهدید محاسبات کوانتومی و رمزنگاری پساکوانتومی}

در بافت روش‌های تولید شبه‌نام، مرور ره‌آواتی آگوستینا و همکاران پیشنهاد می‌کنند که پژوهش‌های آینده تهدید رایانش کوانتومی را پیش‌بینی کنند و هشدار می‌دهند که استفادهٔ فعلی از توابع درهم‌ساز ممکن است برای استفادهٔ بلندمدت مناسب نباشد \cite{rahmawatiagustinaSystematicLiteratureReview2025}. این تنها اشارهٔ مستقیم به تهدید کوانتومی در منابع بررسی‌شده است و جزئیات گذار به رمزنگاری پساکوانتومی \LTRfootnote{\lr{Post-Quantum Cryptography}} — مانند شبه‌نام‌ها و امضاهای مقاوم به کوانتوم مبتنی بر مشبکه‌ها \LTRfootnote{\lr{Lattice-based Cryptography}} — در این منابع پوشش داده نشده است \cite{rahmawatiagustinaSystematicLiteratureReview2025}. بنابراین، آمادگی زیرساخت‌های شبکه اقتضایی خودرویی برای عصر پساکوانتومی، از جمله ارزیابی کارایی الگوریتم‌های پساکوانتومی در محیط‌های با منابع محدود خودرویی و طراحی پروتکل‌های گذار تدریجی، یک مسئلهٔ باز و کار آینده است که مستقیماً از آن یک جملهٔ مرور مذکور پیروی می‌کند \cite{rahmawatiagustinaSystematicLiteratureReview2025}.

\subsection{تحلیل رسمی حریم خصوصی و امنیت}

به کارگیری ابزارهای تحلیل رسمی برای اثبات ویژگی‌های حریم خصوصی و امنیتی پروتکل‌های شبکه اقتضایی خودرویی باید گسترش یابد. مرور ره‌آواتی آگوستینا و همکاران گزارش می‌کنند که ابزار \lr{ProVerif} تنها در دو مطالعه از مطالعات بررسی‌شده به کار رفته و ابزارهایی مانند \lr{ROM}، \lr{ECDLP}، \lr{ECDH}، منطق \lr{BAN}، \lr{AVISPA} و \lr{SPAN} برای تحلیل رسمی حریم خصوصی طراحی نشده‌اند \cite{rahmawatiagustinaSystematicLiteratureReview2025}. تحلیل رسمی می‌تواند آسیب‌پذیری‌های پنهان در طراحی پروتکل‌ها را پیش از اسقرار آشکار کند، اما در منابع بررسی‌شده، حتی طرح‌های رمزنگاری پیشنهادی نیز صرفاً با استدلال توصیفی تحلیل شده و اثبات رسمی — نه کاهش در مدل مهاجم و نه تأیید با ابزار — ارائه نداده‌اند \cite{mariaBAIVEfficientBlockchainBased2022, rahmawatiagustinaSystematicLiteratureReview2025}.

\subsection{مقررات حریم خصوصی و چارچوب‌های یکپارچه}

فقدان چارچوب‌های امنیتی یکپارچه بین‌لایه‌ای و عدم هماهنگی مقررات حریم خصوصی، استقرار جهانی شبکه اقتضایی خودرویی را با مانع مواجه می‌کند. مرور فارسیمدان و همکاران تعامل‌پذیری میان اپراتورها، مراجع جاده‌ای، خودروسازان، شهرداری‌ها و ارائه‌دهندگان سرویس را — با نقش احتمالی دفتر کل توزیع‌شده — از چالش‌های این حوزه می‌نامند \cite{farsimadanReviewSecurityChallenges2025}. از منظر مقرراتی، مقرراتی مانند مقرریت حمایت از داده‌های عمومی اتحادیهٔ اروپا بر کنترل کامل فرد بر اطلاعاتش تأکید می‌کنند و در مه ۲۰۲۴ مقرریت ۲۰۲۴/۱۱۸۳ اتحادیهٔ اروپا چارچوب هویت دیجیتال اروپایی را بنیان‌گذاری کرد \cite{mazzocca2025survey}. استانداردسازی بین‌المللی و هماهنگی مقررات حریم خصوصی از پیش‌نیازهای استقرار جهانی است و همکاری صنعت، دانشگاه و سیاست‌گذاران برای تدوین استانداردهای مشترک ضروری است \cite{farsimadanReviewSecurityChallenges2025, mazzocca2025survey}.
